\documentclass[11pt,a4paper]{article}
\usepackage[T1]{fontenc}
\usepackage[margin=25mm]{geometry}
\usepackage{amsmath,booktabs}

\newcommand{\Dzero}{{}^{5}\mathrm{D}_{0}}
\newcommand{\Fzero}{{}^{7}\mathrm{F}_{0}}
\newcommand{\Ftwo}{{}^{7}\mathrm{F}_{2}}

\title{Interference of different terms for the dipole moments for
  \(\Dzero\) to \(\Fzero\) and \(\Dzero\) to \(\Ftwo\)
  in the \(\mathrm{Eu}^{3+}\) G1 centre}
\author{}
\date{}

\begin{document}
\maketitle

The two lines considered here are from \(\Dzero\) emission in the
\(\mathrm{Eu}^{3+}\) G1 centre in \(\mathrm{CaF}_2\).
In the calculation, the initial eigenstate is level 50, predominantly
\(\Dzero\). Group 1 ends at level 1, predominantly \(\Fzero\);
group 5 ends at level 7, predominantly \(\Ftwo\) with \(M=0\).
Group 5 denotes this particular crystal-field line, not the whole
\(\Ftwo\) multiplet.

\section*{How the amplitudes combine}

Each electric-dipole component \(D_q\), with \(q=-1,0,+1\), is a
coherent sum of contributions. Writing \(D_q^{(\lambda)}\) for the
sum of all terms of tensor rank \(\lambda\), the electric-dipole
strength, apart from a common normalization, is
\[
  S_{\mathrm{ED}} \propto
  \sum_{q=-1}^{+1}\left|\sum_{\lambda=2,4,6}
    D_q^{(\lambda)}\right|^2 .
\]
Thus terms contributing to the same \(q\) can reinforce or cancel
one another; different \(q\) components are squared separately in
this strength. In the two lines considered here, only \(q=0\) is
nonzero, and the amplitudes are real in the logged eigenvector
convention. An overall sign change of an eigenvector flips all terms
for that transition and has no effect on the strength. It is the
relative signs \emph{within} each transition that matter.

For the \(p=0\) intensity parameters used here, the contribution of
one \(A^{\lambda}_{t0}\) to \(D_0\) is, in the PyCF convention,
\[
  D_0[A^{\lambda}_{t0}](i\to f)
  = -10^{10} A^{\lambda}_{t0}(-1)^{\lambda-1}\sqrt{2t+1}
    \begin{pmatrix}\lambda&1&t\\0&0&0\end{pmatrix}
    \langle i|U^{(\lambda)}_0|f\rangle .
\]
The factor \(10^{10}\) converts to the tabulated amplitude units of
\(10^{-10}\,\mathrm{cm}\). For fixed intensity parameters, the same
rank coefficient multiplies the corresponding tensor matrix element
for both transitions.

\section*{Numerical example}

The source is the PyCF run
\texttt{log.20261119\_121153.eucaf2\_o\_matrix\_elements.py.out}.
The eigenstate indices below are one-based.

\begin{center}
\begin{tabular}{lrr}
\toprule
Unit tensor & Group 1: \(\langle 50|U^{(\lambda)}_0|1\rangle\)
            & Group 5: \(\langle 50|U^{(\lambda)}_0|7\rangle\) \\
\midrule
\(U^{(2)}_0\) & \(+6.742141\times10^{-3}\) & \(-2.274218\times10^{-2}\) \\
\(U^{(4)}_0\) & \(-1.542541\times10^{-3}\) & \(+1.447598\times10^{-3}\) \\
\(U^{(6)}_0\) & \(-5.864849\times10^{-4}\) & \(-3.955174\times10^{-4}\) \\
\bottomrule
\end{tabular}
\end{center}

\begin{center}
\begin{tabular}{lrrr}
\toprule
 & & \multicolumn{2}{c}{\(q=0\) amplitude (\(10^{-10}\,\mathrm{cm}\))} \\
\cmidrule(l){3-4}
Parameter & Value used & Group 1: \(50\to1\) & Group 5: \(50\to7\) \\
\midrule
\(A^{2}_{10}\) & \(+7.390000\times10^{-10}\) & \(+0.03151173\)  & \(-0.1062934\) \\
\(A^{2}_{30}\) & \(-3.310000\times10^{-10}\) & \(+0.01728628\)  & \(-0.05830901\) \\
\(A^{4}_{30}\) & \(-4.770000\times10^{-10}\) & \(+0.004905281\) & \(-0.004603362\) \\
\(A^{4}_{50}\) & \(+2.260000\times10^{-10}\) & \(+0.002598418\) & \(-0.002438486\) \\
\(A^{6}_{50}\) & \(-2.737000\times10^{-9}\)  & \(+0.01090525\)  & \(+0.007354351\) \\
\(A^{6}_{70}\) & \(+5.467000\times10^{-9}\)  & \(+0.02352790\)  & \(+0.01586690\) \\
\midrule
Logged total \(D_0\) & & \(+0.09073486\) & \(-0.1484231\) \\
\bottomrule
\end{tabular}
\end{center}

The first four contributions are printed in the log. The two
rank-6 contributions are reconstructed from the logged \(U^{(6)}_0\)
matrix elements and the angular factors above; they reproduce the
logged totals to the precision of the printed inputs. \(A^{2}_{33}\)
was not used in this run. The \(q=\pm1\) electric-dipole components
and all magnetic-dipole components vanish for both lines.

\section*{Why the ordering can reverse}

Grouping the contributions by rank and rounding gives
\begin{align*}
  D_0(50\to1) &\simeq
    \underbrace{0.04880}_{\lambda=2}
    + \underbrace{0.00750}_{\lambda=4}
    + \underbrace{0.03443}_{\lambda=6}
    = +0.09073,\\
  D_0(50\to7) &\simeq
    \underbrace{(-0.16460)}_{\lambda=2}
    + \underbrace{(-0.00704)}_{\lambda=4}
    + \underbrace{0.02322}_{\lambda=6}
    = -0.14842 .
\end{align*}
Rank 6 reinforces ranks 2 and 4 in group 1 but partially cancels
them in group 5. Group 5 still has the larger moment magnitude in
this run because its rank-2 amplitude is much larger.

To see how the ordering can change, hold the eigenvectors and the
rank-4 and rank-6 parameters fixed, and multiply both \(A^{2}_{10}\)
and \(A^{2}_{30}\) by a factor \(s\). In the same amplitude units,
\begin{align*}
  D_0(50\to1;s) &\simeq \phantom{-}0.048798\,s + 0.041937,\\
  D_0(50\to7;s) &\simeq             -0.164602\,s + 0.016179 .
\end{align*}
For example, at \(s=0.4\) the moments are about \(+0.06146\) and
\(-0.04966\), respectively, so group 1 has the larger magnitude.
This is a sensitivity example, not a new fit.

Rank 6 supplies the opposing amplitude in the logged calculation;
it is not a general requirement for cancellation. At fixed
eigenvectors, reversing the sign of the net coefficient of either
rank 2 or rank 4 would make those ranks oppose each other in both groups,
because both \(U^{(2)}_0\) and \(U^{(4)}_0\) change sign between the two
lines in the logged convention. The cancellation would not be equal
in fractional terms: with the logged magnitudes, the rank-4
amplitude is about \(15.4\%\) of rank 2 for group 1 and \(4.3\%\)
for group 5. Rank 6 can instead interfere constructively in one
line and destructively in the other.

The \(\Ftwo,\ M=0\) admixture in the ground state alone therefore
cannot determine the ordering. Each tensor matrix element depends
on the full initial and final eigenvectors, and each line strength
depends on the squared magnitude of the resulting coherent sum.

\end{document}
